% =============================================================================
% Chapter 4 --- 2025 AI Agent Index
% =============================================================================
\chapter{2025 AI Agent Index}
\label{ch:agent-index}

\section{Source document}

\begin{description}
  \item[Title.] The 2025 AI Agent Index --- Documenting Technical and Safety
    Features of Deployed Agentic AI Systems.
  \item[Authors.] Staufer (Cambridge), Feng (Washington), Wei (Harvard
    Law), Bailey (Stanford), Duan (Concordia AI), Yang (UPenn), Ozisik
    (MIT), Casper (MIT, co-senior), Kolt (Hebrew University, co-senior).
  \item[Venue.] \url{https://aiagentindex.mit.edu/}; CC-BY-4.0.
  \item[Pages.] 39.
  \item[Source.] \corpus{source/2025-AI-Agent-Index.pdf}
    (sha256 \texttt{c67a3046\ldots0e267f}).
\end{description}

\section{What it catalogues}

The Index documents 30 state-of-the-art deployed agentic systems across
six documentation categories:

\begin{enumerate}
  \item \textbf{Legal.} Licensing, terms of service, responsible-use
        policies, redistribution rights.
  \item \textbf{Technical capabilities.} Supported modalities, tool
        access, planning horizon, memory.
  \item \textbf{Autonomy and control.} Spectrum from recommender to
        operator; human-in-the-loop requirements.
  \item \textbf{Ecosystem interaction.} External systems the agent
        integrates with (MCP, web browsers, code editors, e-commerce).
  \item \textbf{Evaluation.} Published benchmarks, capability scores,
        safety evaluations, red-team assessments.
  \item \textbf{Safety.} Safety documentation, risk assessments,
        residual-risk statements.
\end{enumerate}

\section{Headline finding for governance}

The Index's top-level empirical finding is directly relevant to our
corpus design:

\begin{quote}\small
\emph{``\ldots we find different transparency levels among agent developers
and observe that most developers share little information about safety,
evaluations, and societal impacts.''}
\end{quote}

Our response is the transparency requirement set extracted in Table~\ref{tab:agentidx-reqs}.

\begin{table}[h]
\centering
\small
\begin{tabularx}{\textwidth}{p{3.2cm}p{1.4cm}p{1.6cm}X}
\toprule
\textbf{Req.\ id} & \textbf{Section} & \textbf{Obligation} & \textbf{Normative text (paraphrased)} \\
\midrule
REG-AGENTIDX-3.6-001 & 3.6 & should & Publish per-agent: risk assessment summary, evaluation results, known limitations. \\
REG-AGENTIDX-3.3-001 & 3.3 & should & Document agent autonomy level and map each state transition to that level. \\
\bottomrule
\end{tabularx}
\caption{Requirements derived from the 2025 AI Agent Index.}
\label{tab:agentidx-reqs}
\end{table}

\section{How we verify the transparency requirement}

Publication is evidenced by this very document plus the structured
records. In practice:

\begin{itemize}
  \item Risk assessment: \corpus{../tb/structured/requirements/tb-business-requirements.json},
        \texttt{riskAssessment} field.
  \item Evaluation results: AI Verify \texttt{process-checklist} plugin
        output plus Moonshot benchmark runs (\cref{ch:tooling}).
  \item Known limitations: the gap entries in this specification's
        coverage tables.
\end{itemize}

This chapter is cited from requirement record
\texttt{REG-AGENTIDX-3.6-001} as
\path{docs/regulations/spec/chapters/04-agent-index.tex}, closing the
loop.
\clearpage
